\documentclass[a4paper,11pt]{article}
\usepackage[margin=24mm]{geometry}
\usepackage{xeCJK}
\setCJKmainfont[BoldFont=IPAexGothic]{IPAexMincho}
\setCJKsansfont{IPAexGothic}
\setCJKmonofont{IPAexGothic}
\usepackage{booktabs,array,multirow,graphicx,amsmath,caption,float}
\usepackage{pgfplots}\pgfplotsset{compat=1.16}
\usepackage{tikz}\usetikzlibrary{positioning,arrows.meta,shapes.geometric}
\usepackage{hyperref}
\usepackage{enumitem}
\setlist{nosep,leftmargin=1.6em}
\captionsetup{font=small,labelfont=bf}
\renewcommand{\figurename}{図}\renewcommand{\tablename}{表}
\renewcommand{\abstractname}{要旨}\renewcommand{\refname}{参考文献}
\newcommand{\IDK}{\texttt{I do not know.}}
\setlength{\parindent}{1em}
\hypersetup{colorlinks=true,linkcolor=black,urlcolor=blue}

\title{\textbf{小学校低学年向け英語会話Bot「Mia」第2世代(MIA2)の開発}\\[1mm]\large 語彙の調査、文脈長1,024の基礎学習、項目(カード)読解、計算とカードの振り分け、実機会話テスト、強化学習の試み}
\author{マーガレットリサーチ\\[1mm]\small 記録: 秘書 いずみ (Claude Sonnet 5.5)}
\date{2026年10月28日}

\begin{document}
\maketitle

\begin{abstract}
本報告は、自作の小型LLM(約4,000万パラメータ)による英語会話Bot「Mia」の第2世代(MIA2、2026年10月23日〜28日)の開発をまとめたものである。第1世代の報告書(2026年10月22日)の後、(a)日常会話の学習不足による不安定さ、(b)計算とカード(知識)への切り替えが働かないこと、(c)カードが1行で情報が少ないこと、が課題として残った。
MIA2では、まず語彙と文脈長を調べ、基礎学習用の語彙は据え置き(読解にだけ出る語は後から足した)、文脈長を256から1,024に伸ばすことを決めた。基礎学習(約11.8億トークン)をやり直し、Qwenで作った長い会話(5.45万会話)で会話SFTを行った。知識は、題名と200〜500トークンの説明文を持つ\textbf{項目}(2,226件、質問追加後で約8.4万問)に作り直した。項目の作り直しと質問の追加などの一連の改良で、50項目931問の自動判定の読解正解率は68\%から80\%になった(1行カードとの同一条件の比較はしていない)。
計算とカードへの切り替えは、現行のMiaを微調整した\textbf{別モデルのルーター}(Plan B)が担当する。ルーターは文章題から式を取り出して電卓に渡し(計算の答え0.976、ルーターv6、凍結テスト)、事実の質問には項目を検索して現行Miaに読ませる(全体の正解率70.4\%、正しい項目を引けた率は下限つきで約73\%)。実機で見つかった弱点は、追加データでルーターv6を作って直した。
評価は、読解、会話、計算、カード、実機会話の5つの基準テストにまとめた。最後に、実機会話の質を強化学習(GRPO型、DeepSeek採点の報酬)で上げようとしたが、2回とも実機会話テストが誤差の範囲で、読解が低下したため採用しなかった。
なお、評価は単一の乱数の種で、採点にLLMを使い、評価問題の一部は開発中の選択にも使っている。値には偏りを含む(第\ref{sec:limits}節)。
\end{abstract}

\section{はじめに}
\subsection{背景と課題}
第1世代のMiaは、語彙を約16,000トークンに絞った自作モデルで、知識カード(1行)を検索して読解モデルに読ませ、計算は電卓に任せる構想であった。その後の実機の会話で、オーナーが次の課題を挙げた。

\begin{enumerate}
\item 普通の会話はそこそこできるが、「あなたはどう?」と聞くと答えられなくなり、会話が不自然になる。TinyStoriesで学習しており、日常会話の学習が不足していると考えた。
\item 計算やカードへの切り替えがうまく働かない。「2+3」には答えるが、「鉛筆2本と鉛筆3本で何本」には切り替わらない。
\item 切り替えてもカードを見て「わからない」になることが多い。カードは1行なので、辞典のように数行にわたって書いた項目にしたほうがよいのではないか。
\item 語彙を増やすのは適当でないと考える。語彙の現状と学年の対応、会話で覚えておくべき長さ(トークン数)を確認したい。
\item 限られた語彙で会話例を大量に作る仕組みが必要である。費用がかかるので、Qwen(手元で動かす)で行う。
\end{enumerate}

\subsection{進め方}
作業は次の順に行い、新しい作業はオーナーの承認を得てから始めた。成果がたまった時点で報告書を書き、作業記録(MIA2.txt)にも追記した。
\begin{center}
\small
語彙・学年・文脈長の調査 $\rightarrow$ 長い会話データの生成 $\rightarrow$ 基礎学習 $\rightarrow$ 会話SFT $\rightarrow$ 項目読解 $\rightarrow$ 計算・カードの振り分け $\rightarrow$ 実機会話テスト $\rightarrow$ 強化学習
\end{center}

\subsection{全体の構成}
最終的なBotは、3つのモデルと1つの検索からなる(図\ref{fig:arch})。ルーターと答え役は、現行のMiaをコピーして微調整した\textbf{別のモデル}である。現行のMiaは、会話と、項目を読んで答える役を兼ねる。

\begin{figure}[H]
\centering
\begin{tikzpicture}[node distance=7mm and 10mm, box/.style={draw,rounded corners,minimum height=8mm,align=center,font=\small}, >=Stealth]
\node[box] (u) {ユーザーの発言};
\node[box,right=of u] (r) {ルーターMia v6\\(calc / card / chat)};
\node[box,above right=2mm and 14mm of r] (calc) {電卓};
\node[box,right=of calc] (ans) {答え役Mia\\(結果を文にする)};
\node[box,right=14mm of r] (ret) {項目の検索\\(Qwen3-Embedding-0.6B\\+採用の下限)};
\node[box,below right=2mm and 14mm of r] (chat) {現行Mia\\(会話)};
\node[box,right=of ret] (read) {現行Mia\\(項目を読む)};
\draw[->] (u)--(r);
\draw[->] (r)--(calc) node[midway,above,font=\scriptsize,sloped]{calc 式};
\draw[->] (calc)--(ans);
\draw[->] (r)--(ret) node[midway,above,font=\scriptsize]{card 題名};
\draw[->] (ret)--(read);
\draw[->] (r)--(chat) node[midway,below,font=\scriptsize,sloped]{chat};
\end{tikzpicture}
\caption{Botの構成。ルーターが3つの経路に振り分ける。}\label{fig:arch}
\end{figure}

\section{語彙・学年・文脈長の調査(10月23日)}
\subsection{語彙}
トークンではなく見出し語(活用をまとめた数、WordNetと頻度で簡易に判定)で数えた。現在のトークナイザーは15,946トークンで、うち単語として実在する13,206トークンが、見出し語7,899語にあたる。語彙を「約15,900」と言うより、人間の数え方では約7,900語を1トークンで持っている。

学習データで実際に使われる語は、会話データ(Qwen生成30万会話、延べ1,304万語)で見出し語9,958語、TinyStories(先頭2,000万トークン)で9,325語であった。使用頻度は偏っており、会話データでは上位500語が延べの88.7\%、上位2,000語が98.6\%を占める。\textbf{延べ語数の99.6\%は、すでに1トークンで持っている語}で、持たない語は見出し語で約4,200あるが、延べでは0.4\%(名前、擬音など)である。この結果から、語彙は増やさない(増やしても得るものが小さい)と決めた。

\subsection{学年の目安}
物差しにはDale-Chall(新版2,884語、小学4年生の8割が知っている語)\cite{dalechall}を使った。会話データの延べ語数の94.1\%がDale-Challに入り、TinyStoriesは90.5\%である。見出し語の数では25\%にとどまるが、残りは固有名詞とまれな語で、延べでは数\%にすぎない。使われる語の中心は、\textbf{小学校低〜中学年(2〜4年生)}の語彙と推定する。年齢別習得表(Kupermanなど)は入手できず、学年はDale-Challと頻度からの推定である。

\subsection{文脈長}
当時の文脈長(block\_size)は256トークンで、会話データの平均は54トークン、最長246トークン(最長4往復)であった。実測では、1発言が平均16.3トークン、1往復41、4往復110で、1,024トークンを会話で埋めるには約30往復が要る。文脈長を変えて増えるパラメータは位置埋め込みだけで、256から1,024では約0.4M(全体の約1\%)である。ただし位置埋め込みを広げるため、基礎学習のやり直しが必要になる。オーナーの希望により、文脈長を\textbf{1,024}にすることを決めた。

\section{データの作成}
\subsection{長い会話(Qwen)}
1,024トークンに合う長い会話は公開データにほとんどなく(10往復以上のものは見つからなかった)、Qwen3.8-27B(INT4)を手元のGPU(RTX 4090)でvLLMにより動かして作った。API費用はかからず、時間と電気代だけである。試作(10往復と30往復を各100会話)では、Qwenは往復数を数えられず頼んだ数より少なく書く(30往復を頼んで平均27往復)こと、話がよく脱線することが分かった。小学4年生の語の割合は92.5\%(既存データは94.1\%)であった。

本番生成は、10往復・20往復・30往復を3:2:1の割合で、約13時間49分かけて行った(表\ref{tab:long})。

\begin{table}[H]\centering\small
\caption{長い会話の生成結果(合計54,525会話、2,206万トークン)}\label{tab:long}
\begin{tabular}{lrrrr}\toprule
頼んだ往復数 & 会話数 & 総トークン(万) & 実際の往復(平均) & 1会話の平均トークン\\\midrule
10往復 & 29,141 & 803 & 8.3 & 275\\
20往復 & 17,912 & 888 & 17.6 & 496\\
30往復 & 7,472 & 515 & 26.5 & 689\\\bottomrule
\end{tabular}\end{table}

不採用は約4.9万件で、読めない形式3,694、ターン数が奇数3,575、1,024超え801、評価の固定問題との重複131などである。サンプルを読んだ範囲では、話は最後まで続き語彙もやさしいが、Miaの返事が毎回「共感+質問」に偏る傾向があった。

\subsection{基礎学習用のお話}
基礎学習には、TinyStories\cite{tinystories}(約5.49億トークン)に加え、SimpleStories\cite{simplestories}(MIT、211.6万本、6.29億トークン)を使うことにした。SimpleStoriesの調査(2万本の無作為抽出)では、平均310トークン、最長900、1,024超は0本で、小学4年生の語の割合は延べ89.2\%、「Once upon a time」を含む話は0.6\%であった(論文はTinyStoriesで59\%と指摘)。合計約11.8億トークンである。

\section{基礎学習(文脈長1,024)}
モデルは10層・8ヘッド・512次元(40.17M)で、単語埋め込みはQwen3-Embedding\cite{qwen3emb}を主成分分析で512次元にしたものを固定する。事前学習の量は、約20トークン/パラメータという目安\cite{chinchilla}を下限の参考にし、実際は約35億トークン(約88トークン/パラメータ、3エポック)を学習した。データの重複なしの量(11.8億トークン)でも約29である。

学習は、各エポックで485万話をシャッフルし、1話を1行(\texttt{<bos>}から\texttt{<eos>})として話をまたがず、長さの近い話をそろえてバッチにする。3エポック(120,789ステップ、約35億トークン)、lr $10^{-3}$、cosineである。
最初の試行では、話が最長で約950トークンのため、全ての話を位置0から始めると700〜1,023番の位置埋め込みがほとんど学習されない見込みであった(実測ではなく見込みである)。そこでオーナーの提案により、\textbf{50\%の話を位置0から、残り50\%を$0$〜$1024-$(話の長さ)の乱数の位置から}始めるやり直しを行った。検証損失を「位置0から」[start]と「1,024の末尾にそろえる」[end]の両方で測ると、差は0.004以下で、後ろの位置も同じ精度で学習できた(表\ref{tab:pre})。

\begin{table}[H]\centering\small
\caption{基礎学習の最終の検証損失(step 120,000、各3,000話)}\label{tab:pre}
\begin{tabular}{lrr}\toprule
データ & [start] & [end]\\\midrule
TinyStories & 1.0981 & 1.1022\\
SimpleStories & 1.5099 & 1.5133\\\bottomrule
\end{tabular}\end{table}

予定の120,789ステップのうち約700ステップ(0.6\%)を残して、実装のバグ(エポックごとの捨てるバッチ数が変わる)で停止した。学習率は116,000以降下限で損失も横ばいだったため、実害は小さいと判断し、step 120,000のチェックポイントを使った。所要は約3時間である。

\section{会話SFT}
短い会話326,115と長い会話53,406の合計379,521会話(3,794万トークン、損失の対象2,123万)で教師あり微調整を行った。3エポック、lr $3\times10^{-4}$、埋め込み固定、文脈長1,024で、約9分22秒である。検証損失は開始前3.9424から1.6415(短い1.6285、長い1.6472)になった。5エポックにすると、エポック4から5で短い会話の検証損失が上がり(1.6502から1.6788)、過学習の兆しが出たため、3エポックで足りると判断した。

オーナーの試し会話(約400トークン)では、AIであることを正直に言う、滑り台やニューヨークの話題が自然、という利点があった。一方、綴り間違いによる取り違え、同じ会話内の矛盾(「My cat is named Coco」と言った直後に猫の名前を聞くと「I do not know your cat's name」と答えた。3エポック版と5エポック版の両方)、代名詞の取り違え(「My son lives there」に「He is a famous place.」)、意味の通らない文が見られた。400トークンで起きているため、文脈の長さではなく、モデルの理解力とデータの偏りが原因と見ている。

\section{項目(カード)による読解}
\subsection{項目の設計}
カードは短文では情報が足りないとのオーナーの指摘を受け、\textbf{題名+200〜500トークンの説明文+質問}の「項目」にした。小学校低学年(1〜2年)がわかる言葉で、1文12語以内とする。質問は1項目20問(wh 9、yes 2、no 2、count 3、unk 4)で、unkは本文に書かれていないことを聞き、答えは必ず\IDK である。

題材は、12のテーマ(動物・昆虫・植物・地球・天気・宇宙・からだ・食べ物・乗り物・家・学校と仕事・遊び)ごとにQwenに挙げさせ、語彙の条件とQwenの判定で整えて1,155題材とし、観点を増やして1,076題材を追加し、合計2,231題材とした。

本文と質問の作り方は2つ試した。Qwenが書いた本文をDeepSeekが事実確認する方法は、間違いの見逃しと新しい誤りがあり使えなかった。採用したのは、\textbf{DeepSeekが本文と質問20問を1回で直接作り、Qwenが各質問と答えを検証する}方法である(失敗は3回までやり直し)。結果は2,226項目(約4.4万問)、DeepSeek費用は約\$4.95(見積もり)であった。Qwenの検証のok率は、最初の1,151項目で95.9\%であった。

\subsection{語彙の追加}
読解データにだけ出る割れる語(基礎データと会話データに片の並びが一度も出ない語)は751語あった。出現5回以上などで絞った145語(nougat, bonsai, caribouなど)と中間の断片121個を足し、語彙を15,946から\textbf{16,212}にした(新語彙ts\_v4)。基礎モデルは学習し直さず、新語の埋め込みは旧語と同じ手順で計算し、旧語の再計算が保存値と一致する(cos 1.0)ことを確認してから行を足した。語彙を増やさない方針の例外であり、オーナーが「読解にだけ出る語なら基礎からやり直さなくてよい」と提案した。

\subsection{読解SFTと基準テストの確定}
会話SFT済みのモデルから読解SFTを行った。1項目につき質問5問の多ターン会話を複数サンプルとし、会話の復習を混ぜた。質問は、同じ項目の本文に対してDeepSeekが既存の質問と重ならない新しい質問を作り(39,934問、費用の見積もり約\$3.85)、noの弱さに合わせて配分を変えて倍にした(第5版)。

評価は、学習に使わない50項目の全931問を貪欲生成で答えさせた、自動判定の正解率である。unkは答えが\IDK と完全一致、他の型は模範解答の内容語(3文字以上)の6割以上が答えに出ていて、かつ\IDK でないこと(yes/noは先頭語も一致)を条件にする。判定は内容語の一致なので、人が見るより甘い。

\begin{table}[H]\centering\small
\caption{読解SFTの推移(読解テスト931問の正解率\%と会話テスト)}\label{tab:read}
\begin{tabular}{lrrrrrrr}\toprule
モデル & wh & yes & no & count & unk & 全体 & 会話\\\midrule
会話SFTのみ & - & - & - & - & - & 7 & 85.0\\
読解第1回(旧語彙、1,151項目) & 62 & 70 & 57 & 70 & 86 & 68 & -\\
読解第3回(2,226項目) & 71 & 87 & 54 & 82 & 83 & 75 & 84.1\\
\textbf{読解第5版(質問を倍、確定)} & 78 & 76 & 71 & 81 & 90 & \textbf{80} & \textbf{84.6}\\\bottomrule
\end{tabular}\end{table}

第5版は、wh・no・unkが大きく改善した一方、yesは87から76に下がった(追加の質問がnoに偏った影響の可能性)。会話の検証損失は1.64から1.79に悪化したが、会話テストの点は保たれた。2026年10月26日に、この第5版を\textbf{現行モデル}として確定した。

\section{会話テスト}
オーナーの提案により、5回のやりとりに点をつけるテストを作った。会話SFTの検証会話から、短い50本と長い50本(最初の5発言)を選んで凍結し、ユーザー発言を固定してMiaに貪欲生成で答えさせる。DeepSeekが1回20点(答える8・覚えている6・ことば4・きれいさ2)で採点し、1回の平均$\times5$を100点満点とする。基準は\textbf{標準的な小学校低学年(7歳)の子どもが約90点}(1回18点)で、質問を無視した返事は最大10点、矛盾・取り違えは最大12点、不適切な助言は最大8点とする。最初の基準では89.0点と甘すぎたため、2回直して厳しくした。採点者は甘めの傾向があるので、絶対値より比較に使う。結果は、会話SFTのみ85.0点、読解SFT後(確定)84.6点である。

以後、読解テスト(約1分、費用なし)と会話テスト(約2分、約\$0.07)を\textbf{基準テスト}とし、問題と採点基準は変えず、変えるときは旧版と並べて記録する決まりにした。

\section{計算の振り分け:ルーターMia(Plan B)}
\subsection{課題と方針}
従来の計算は、ルーターの頭が計算と判断したら正規表現で式を取り出す一方通行で、文章題から式を作れず、返事も固定文であった。また、答える側のMiaは電卓の結果を読めない(\texttt{Facts: 12 - 1 = 11}を渡して正解5\%)。

判断と答えを同じモデルにやらせる案(Plan A)は、小さなモデルへの負担と会話の劣化が心配であった。そこで、オーナーの提案により\textbf{Plan B}、すなわち現行Miaを微調整した\textbf{別のルーターMia}(判断と式づくり)と、別に学習した\textbf{答え役Mia}(結果を自然な文にする)の構成にした。ルーターの出力は特殊トークンなしの平文で、\texttt{calc 2 + 3}、\texttt{card 題名}、\texttt{chat}である。

\subsection{データ}
文章題を14の型(増える、演算語なし、のこり、差、かけ算、わり算、2段、関係ない数字入り、3項、おつり、年上年下、半分2倍、式だけ、続きの質問)で、テンプレートから作った。大きい数(最大約3,000)を3割にした。式と答えは電卓で検算する。言い換えはQwenが3通り作り、\textbf{数の集まりが元と同じで、Qwenが独立に解いて元の答えと一致するものだけ}を採用した(最初は関係ない数字が足し算に混ざる言い換えがあり、この検査で落とした)。数字を含むが計算ではない発言(「I am 7 years old.」)も、chatの否定例として加えた。

\subsection{経過と結果}
\begin{table}[H]\centering\small
\caption{ルーターの経過(計算:判断/式/答え、会話:誤呼び出し率)}\label{tab:router}
\begin{tabular}{lll}\toprule
版 & 計算(判断/式/答え) & 会話の誤呼び出し\\\midrule
v1(単発の問題だけ) & 1.000 / 0.968 / 0.987(旧テスト) & 0\\
v2(型と大きい数を追加) & 0.978 / 0.627 / 0.662(履歴つきテスト) & 0.026\\
\textbf{v3(履歴に前の計算・会話)} & \textbf{1.000 / 0.943 / 0.986} & \textbf{0}\\\bottomrule
\end{tabular}\end{table}

v1は凍結テストで高く出たが、手書き20問で「2 plus 3」や3項、4桁の数が崩れることが分かった。型を足したv2は、会話の途中で履歴に前の計算があると前の問題の数字を混ぜる欠陥が出た。v3は、履歴に前の計算や会話をつけて学習し、解消した。

答え側は、現行Miaを別コピーで微調整した\texttt{answer\_mia\_v1}が答えの正しさ0.991(800件)、固定文に落ちる率0.005であった。現行Miaに復習4万件を混ぜて1つにまとめる案(\texttt{answer\_mix\_v1})は、答え0.981、読解79\%、会話83.1と、少し下がった。オーナーは、当面は別コピーのまま保持すると決めた。

\section{カードへの拡張}
\subsection{ルーターへの追加}
ルーターの出力に\texttt{card 題名}を足し、事実の質問を項目の検索に回す。学習データは、項目の質問のうち題名を含むもの(12,224)、「Tell me about X」型(4,152)、代名詞の質問(2,177)、題名を含むが事実を聞かない発言(chatへ、3,869)、前の話題を引きずらない切り替え(2,443)、項目にない話題の質問(3,010)などである。評価用に150題名(読解テストの50項目と無作為100)を学習から外した。

\subsection{検索と採用の下限}
検索は、項目の本文をQwen3-Embedding-0.6B\cite{qwen3emb}で埋め込み、質問との近さに、ルーターが出した題名の手がかりを小さく加える。採用は2段の下限で決める。質問との類似度が上の下限(0.70)以上なら採用し、下の下限(0.50)以上で上の下限未満のときは、ルーターが出した話題の語が項目の本文に出る場合だけ採用する。どちらでもなければ「項目なし」とし、\IDK と答える。Miaの単語埋め込み(\texttt{wte})の平均で検索する案も試したが、正しい項目を1位に引けた率が62\%(語の重みなしでは27\%)で、同じ条件(下限なし、1位)のQwen3-Embeddingの77.3\%に及ばなかった。

\subsection{全体テストと下限の選択}
読解テストの50項目931問を、ルーター、検索、現行Miaの読解の順に通した(題名はルーターの学習に使っていない)。項目を最初から渡した場合は80\%である。

\begin{table}[H]\centering\small
\caption{カードの全体テスト(931問)と、項目にない話題215問での\IDK の率}\label{tab:card}
\begin{tabular}{@{}l@{\hspace{3mm}}r@{\hspace{3mm}}r@{\hspace{3mm}}r@{}}\toprule
構成 & 全体の正解率 & 項目を引けた率 & 項目なしの\IDK\\\midrule
項目を最初から渡す(参考) & 80.0\% & - & -\\
題名索引のみ & 70.0\% & 67.6\% & -\\
埋め込み検索(上位1枚、下限なし) & 74.8\% & 77.3\% & -\\
上位3枚を順に読む & 73.6\% & 77.3\% & -\\
下限 0.40/0.60 & 73.6\% & 76.0\% & 57.7\%\\
\textbf{下限 0.50/0.70(採用)} & \textbf{70.8\%} & 72.7\% & \textbf{67.0\%}\\
下限 0.60/0.80 & 63.8\% & 63.4\% & 83.3\%\\\bottomrule
\end{tabular}\end{table}

下限を上げると、作り話は減るが、正しい項目の質問も捨てられる(取りこぼしと作り話の交換)。オーナーは、バランスのよい0.50/0.70を採用した。上位3枚を順に読む方式は、別の項目から答えを作ってunkの正解率を落とす(89\%から79\%)ため採用しなかった。外れの主因は、質問が題名でなく部分の名前だけを言う場合(cornの項目の「What covers the cob?」)と、項目の重なり(honeydewとaphid)である。

\subsection{ルーターv6:実機の弱点を直す}
オーナーの実機テストで、雑談が続くと整った質問でもカードを引かない、「You know what is Sun?」「Tell me more about Moon」などのくだけた言い方や、「grass hopper」のような分けた題名を引かない、という弱点が見つかった。項目はすべてあったので、原因はルーターにあった。雑談の履歴つきの事実の質問(7,000)、くだけた言い方(6,228)、同じ話題の好みの後の質問(1,500)、好み・気持ちのchat(6,500)、分けた題名(506)をテンプレートで足し(21,734件)、v6を作った(表\ref{tab:v6})。

\begin{table}[H]\centering\small
\caption{ルーターv5とv6}\label{tab:v6}
\begin{tabular}{lrr}\toprule
テスト & v5 & v6\\\midrule
新テスト(学習と別の言い方・別の題名)のカード判断 & 0.610 & \textbf{0.923}\\
新テストの題名一致 & 0.488 & 0.778\\
凍結テストのカード判断 & 0.995 & 0.993\\
凍結テストの計算の答え & 0.977 & 0.976\\
会話の誤呼び出し率 & 0.0013 & 0.0003\\
全体テスト(下限0.50/0.70) & 70.8\% & 70.4\%\\
項目なし話題の\IDK & 67.0\% & 65.6\%\\\bottomrule
\end{tabular}\end{table}

新しい弱点だけが大きく伸び、凍結テストと計算はほぼ維持された。全体テストはわずかに下がり(70.8から70.4\%、項目なしの断りは67.0から65.6\%)、これは誤差程度とみている。実機でうまくいかなかった発話は、すべて項目を引くようになり、「You like Moon?」は意図どおりchatのままであった。

ただし、新テスト(0.610から0.923)は学習と同じテンプレート生成器(言い回しと題名だけを分けたもの)で作っており、同じ分布での改善である。また、実機会話テストの6場面はv6を作る動機になったオーナーの実機ログで、学習用の言い回しもその失敗例を模している。テストの場面そのものは学習に使っていないが、場面の型は開発の標的であった。

\section{実機会話テスト}
\subsection{テストの設計}
基準テストでは測れない不自然さ(広い依頼への答え、代名詞の続き、雑談からカードへの出入り、文の崩れ、事実の間違い)を測るため、オーナーの指示で5つ目の基準テストを作った。ボット全体を通す固定の複数ターンの場面24個(計120発言)で、オーナーの実機ログ6場面のほか、計算、カード、雑談、記憶、崩れた英語、話題の切り替えを含む。ユーザー発言は固定で、場面ごとに乱数を固定する。DeepSeekが会話テストと同じ20点式で採点し、基準は7歳の子どもで約90点、\textbf{合格ラインは80点}とする。事実や数の間違いは最大8点に抑える。学習の報酬や学習データには使わない決まりにした。

\subsection{結果}
現行の構成(ルーターv6、現行Mia、答え役、下限0.50/0.70)の点は\textbf{77.9〜78.2点}で、合格ラインに届かなかった。2回の測定は、77.9が120発言、78.2が110発言(1場面を採点できなかった)の平均である。2回はボットの返事が同じ110発言で、採点だけが0.3点違い、これが採点の再現性の目安になる。以下の内訳は77.9の回の値で、120発言のうち80点以上は87、60点未満は26、間違い・作り話は15であった。場面別では、記憶・雑談・大きい数の計算・天気・話題の切り替えが90点で、項目にない話題(55点)、面積(62点)、お金の続きの計算(65点)、accordion(66点)が低かった。

\begin{table}[H]\centering\small
\caption{実機会話テストで80点未満だった発言(現行構成、経路別)}\label{tab:live}
\begin{tabular}{lrr}\toprule
経路 & 80点未満 & 発言数\\\midrule
会話 & 7 & 63\\
\textbf{カード} & \textbf{22} & 42\\
計算 & 4 & 15\\\bottomrule
\end{tabular}\end{table}

失点の半分以上がカードの経路にある。原因は、(1)項目にない話題で似た項目から作り話をする(「What is a bridge?」に、roadの項目から「A bridge is a long hole in the ground.」)、(2)広い依頼で読む役が外す(「Tell me about Mars.」に「The moon is a planet.」、「Is Mars a moon?」に「Yes」)、(3)代名詞の続き(「tell me more about it」)で検索が項目を選ばない、(4)計算の型の不足(面積が\texttt{5 + 2}になる、お金の続きで前の残りを使わない)、である。

\section{強化学習の試み}\label{sec:rl}
実機会話は「正解がなく、揺らぎの中から良いものへ誘導したい」とのオーナーの考えから、強化学習(GRPO型\cite{grpo})を試した。現行Miaの別コピーを方策とし、新しく作った場面(実機会話テストと別の題材)の途中まで進めて、同じ状況の返事を8本サンプルする。DeepSeekが8本を実機会話テストと同じ20点式で採点し(カードの回は項目の本文を、事実の誤りの確認に使う)、群内で標準化した報酬を利得にして方策勾配を1回更新する。元のモデルとのKLの制限も加えた。1ステップは8場面で、100ステップの費用はDeepSeekの使用量の見積もりで約\$1.0、約14分である。訓練中の報酬は、温度1.0のサンプルを別の採点プロンプト(項目の本文つき)で採点した値で、実機会話テストの点とは設定が違うため、数値を直接比べない。

\begin{table}[H]\centering\small
\caption{強化学習の2回の結果(100ステップ)。すべて単一の乱数の種}\label{tab:rl}
\begin{tabular}{lrrr}\toprule
モデル & 実機会話テスト & 読解テスト & 訓練中の報酬(20点満点)\\\midrule
現行Mia & 77.9 & 80\% & -\\
第1回 50ステップ(復習なし、$\beta=0.02$) & 77.6$^*$ & 76\% & 11.7$\rightarrow$13.7\\
第1回 100ステップ & 79.0 & \textbf{71\%} & (同上)\\
第2回 50ステップ(復習+$\beta=0.1$) & 78.0 & 78\% & 11.7$\rightarrow$12.3$^\dagger$\\
第2回 100ステップ & 78.4$^*$ & 77\% & (同上)\\\bottomrule
\end{tabular}\\[1mm]
\parbox{0.88\linewidth}{\scriptsize $^*$1場面を採点できず110発言の平均。共通の110発言で比べ直すと、現行Mia 77.86、第1回50 77.59、第1回100 78.41、第2回50 77.45、第2回100 78.36で、差は$-0.4$〜$+0.6$点。 $^\dagger$最初の20ステップの平均から最後の10ステップの平均。途中の10ステップ移動平均の最大は13.4。第1回は最初と最後の各20ステップの平均。}
\end{table}

訓練中の報酬は上がったが、実機会話テストの差(最大$+1.1$点、共通の発言で比べると$-0.4$〜$+0.6$点)は誤差の範囲である。読解は、復習なしでは71\%、復習を混ぜても77〜78\%に下がり、特にnoの質問が弱くなった(71\%から51〜58\%)。基準を保てないため、\textbf{2回とも採用しなかった}。会話の経路は失点が少なく、失点の主因が返事の文ではなく、ルーター・検索・計算の型にあるため、返事を強化する学習では伸びにくかったと考える。

\section{考察}\label{sec:limits}
\subsection{うまくいったこと}
\begin{itemize}
\item \textbf{項目への作り直し}。カードを辞典のような項目にして読解を学習し、項目の作り直しに質問の倍増と語彙の追加を含む一連の改良で、読解は68\%から80\%になった(第1回は旧語彙・1,151項目、第5版は質問倍増で条件が複数違い、1行カードとの同一条件の比較はしていない)。なお、オーナーの試用では、基盤のMiaが項目の内容をかなり覚え、カードなしでも答えることがあるという印象であったが、これは本報告では測定していない(測るなら、読解テストの50項目の題名だけを聞く、などの評価が要る)。
\item \textbf{ルーターの別モデル化}。判断と式づくりを別モデルにしたことで、現行Miaの読解80\%と会話84.6点を保ったまま、計算の答え0.976〜0.977(版・テストにより異なる)を得た。
\item \textbf{検査つきの合成データ}。言い換えを独立に解かせて答えを確かめる、数の集まりが同じかを確かめる、という検査で、意味が変わったデータを除けた。
\item \textbf{実機で見つけて直す循環}。基準テストは決めた範囲の後退を見つけるのに強く、実機の試用は範囲の外の弱点を見つける。両方を使うことで、ルーターv6の弱点を見つけて直せた。
\end{itemize}

\subsection{限界と注意}
\begin{itemize}
\item すべての結果は\textbf{単一の乱数の種}の値で、実機会話テストの数点の差は誤差である(120発言)。同じ返事の再採点でも約0.3点違った。会話テストの差(85.0、84.6、84.1)も同程度で、「保たれた」とは言えても「改善した」とは言えない。
\item 文脈長1,024の効果は、256との比較や語彙据え置きとの切り分けでは測っていない。Botは、ルーターも会話も直近3往復(約80〜110トークン)しか渡しておらず、カードの経路は質問だけで読む。会話テストも最初の5発言までで、長い会話の一貫性の評価は未実施である。代名詞の続きに弱いこととも関係すると考える。
\item 会話テストと実機会話テストの採点は\textbf{LLM(DeepSeek)}で、甘めの傾向があり、絶対値は比較用である。読解の自動判定は内容語の一致なので、文が崩れていても正解になりうる。
\item 評価の50項目などは開発中の選択(エポック数、下限など)にも使っており、\textbf{選択の偏り}を含む。ルーターの凍結テストは、学習と同じ生成器(テンプレート+Qwen言い換え、別シード)で作ったため、実際の言い回しより易しい。
\item 項目の内容はDeepSeekが作り、Qwenが検証したもので、事実の誤りが残る(「No, a tomato is not a fruit.」は項目の記述どおりだが、正しくは果物である)。
\item 答え役の訓練データは、検査器のカンマ入りの数の不具合で、1,000以上の答えの文が生成側で全部落ちていた(直したが再生成は未実施)。答えテスト0.991は、1,000以上の答えの訓練が薄いままの値である。前提のない質問(「7+7と計算」)、訂正、複数の事実をまたぐ計算も未対応である。
\item DeepSeekの費用は、プログラムが使用量から見積もった値である(項目\$4.95、質問追加\$3.85、会話テスト約\$0.07/回、実機会話テスト約\$0.03/回、強化学習約\$1.0/回)。
\item ライセンスは2026年10月28日にHugging Faceのカードと規約で確認した。TinyStoriesはCDLA-Sharing-1.0(データを再配布する場合のみ同じ条件が付き、学習したモデルの公開は妨げない)、SimpleStoriesはMIT、Qwen3-EmbeddingはApache-2.0、DeepSeekのAPI規約は出力を他のモデルの学習に使うことを認めている。会話データの生成に使ったQwen(RedHatAI/Qwen3.8-27B-INT4)もApache-2.0で、カードに出力の利用制限の記載はない。以上は規約の読み取りであり、法的な助言ではない。
\end{itemize}

\section{まとめと今後}
MIA2では、基礎学習用の語彙は据え置いて(読解にだけ出る語を後から足し16,212語)文脈長を1,024にし、基礎学習・会話SFT・項目読解を作り直した。計算とカードへの切り替えは、現行Miaを変えずに、別のルーターで実現した。現在のBotは、会話はスムーズに受け答えでき、計算は文章題から式を取り出して電卓で答え、事実の質問は項目を引いて読んで答える。5つの基準テストの値は、読解80\%、会話84.6点、計算の答え0.976(ルーターv6、凍結テスト)、カードの全体70.4\%(項目にない話題の断り65.6\%)、実機会話77.9点(合格ライン80)である。

未実施の案は次のとおりである。
\begin{enumerate}
\item ルーターが出した題名が項目の題名と完全に一致するなら、下限を通さず採用する(代名詞の続きの対策、学習不要)。
\item 読解側の追加学習(別コピー)。広い依頼、「Do you know\dots」、本文にない関係のyes/noを断る例を足す。
\item ルーターv7。前の計算結果を使う型、面積などの掛け算の型を足す。
\item 部分の名前の索引。
\end{enumerate}
強化学習は、失点の原因が減ってから、もう一度試す価値がある。

\begin{thebibliography}{99}
\bibitem{tinystories} R.~Eldan and Y.~Li, ``TinyStories: How Small Can Language Models Be and Still Speak Coherent English?'' arXiv:2305.07759, 2023.
\bibitem{simplestories} L.~Finke, C.~Sreedhara, T.~Dooms, M.~Allen, E.~Zhang, J.~D.~Rodriguez, N.~Nabeshima, T.~Marshall, and D.~Braun, ``Parameterized Synthetic Text Generation with SimpleStories,'' arXiv:2504.09184, 2025. データ: Hugging Face \texttt{SimpleStories/SimpleStories}.
\bibitem{chinchilla} J.~Hoffmann et al., ``Training Compute-Optimal Large Language Models,'' arXiv:2203.15556, 2022.
\bibitem{qwen3emb} Y.~Zhang et al., ``Qwen3 Embedding: Advancing Text Embedding and Reranking Through Foundation Models,'' arXiv:2506.05176, 2025.
\bibitem{dalechall} J.~S.~Chall and E.~Dale, \textit{Readability Revisited: The New Dale-Chall Readability Formula}, Brookline Books, 1995.
\bibitem{grpo} Z.~Shao et al., ``DeepSeekMath: Pushing the Limits of Mathematical Reasoning in Open Language Models,'' arXiv:2402.03300, 2024.
\end{thebibliography}
\end{document}
